\subsection{COMPLETE SPACES}

In a uniform space X, a Cauchy filter need not have a limit point.

Examples. 1) Consider the sequence \((u_{n})\) on the rational line \(\mathbf{Q}\) defined

\[
\text { by } \quad u _ {n} = \sum_ {p = 0} 2 ^ {- p (p + 1) / 2}. \quad \text { If } \quad m > n \quad \text { we   have }\tag{1}
\]

\[
\left| u _ {m} - u _ {n} \right| \leqslant 2 ^ {- n (n + 3) / 2}
\]

and therefore \((u_{n})\) is a Cauchy sequence. But this sequence has no limit in \(\mathbf{Q}\) ; for if the rational number \(a / b\) were a limit of \((u_{n})\) , then by (1) we should have for all \(n\)

\[
\left| a / b - h _ {n} / 2 ^ {n (n + 1) / 2} \right| \leqslant 1 / 2 ^ {n (n + 3) / 2}
\]

where \(h_n\) is an integer (depending on \(n\) ); that is,

\[
\left| a. 2 ^ {n (n + 1) / 2} - b h _ {n} \right| \leqslant b. 2 ^ {- n}
\]

for all n.~Now the left-hand side of this inequality is an integer for all n, and must therefore be zero whenever n is greater than an integer \(n_{0}\) such that \(b < 2^{n_{0}}\) ; we should therefore have \(a/b = u_{n}\) for all \(n > n_{0}\) , which is absurd.

\begin{enumerate}
\def\labelenumi{\arabic{enumi})}
\setcounter{enumi}{1}
\tightlist
\item
  Let X be an infinite set, and consider the uniformity of finite partitions on X (§ 2, no. 2, Remark 1). Every ultrafilter F on X is a
\end{enumerate}

Cauchy filter with respect to this uniformity. For if \((\mathbf{A}_{i})\) is a finite partition of X and

\[
\mathbf {V} = \bigcup_ {i} (\mathbf {A} _ {i} \times \mathbf {A} _ {i})
\]

the corresponding entourage, then at least one of the \(A_{i}\) belongs to F (Chapter I, § 6, no. 4, Corollary to Proposition 5), and \(A_{i}\) is V-small. On the other hand, X is an infinite discrete space, hence is not compact, and consequently there are ultrafilters on X which do not converge.

DEFINITION 3. A complete space is a uniform space in which every Cauchy filter converges.

In a complete space every Cauchy sequence (no. 1) is therefore convergent. Example. On a discrete uniform space X a Cauchy filter is a trivial ultrafilter (Chapter I, § 6, no. 4), hence convergent; consequently, X is complete.

From Definitions 2 and 3 of no. 1 and Proposition 2 of no. 1, we deduce the following proposition, known as Cauchy's criterion:

PROPOSITION 6. Let \(\mathfrak{F}\) be a filter on a set X, and let \(f\) be a mapping of X into a complete uniform space \(\mathbf{X}'\) . Then \(f\) has a limit with respect to \(\mathfrak{F}\) if and only if the image of \(\mathfrak{F}\) under \(f\) is a Cauchy filter base.

This criterion shows the importance of complete spaces in all questions involving the notion of limit: if a function takes its values in a complete space we can prove the existence of a limit without knowing in advance the value of the limit; this would be impossible if the definition of limit were the only criterion of convergence at our disposal.

A uniformity which is finer than the uniformity of a complete space need not be a uniformity of a complete space (Exercise 2). However, we have the following proposition:

PROPOSITION 7. Let \(\mathcal{U}_1, \mathcal{U}_2\) be two uniformities on a set \(X\) , and let \(\mathfrak{C}_1, \mathfrak{C}_2\) be the topologies induced by these uniformities respectively. Suppose that \(\mathcal{U}_1\) is finer than \(\mathcal{U}_2\) , and that there is a fundamental system of entourages for \(\mathcal{U}_1\) which are closed in \(X \times X\) in the topology \(\mathfrak{C}_2 \times \mathfrak{C}_2\) . Then a filter \(\mathfrak{F}\) on \(X\) converges in the topology \(\mathfrak{C}_1\) if and only if it is a Cauchy filter in the uniformity \(\mathcal{U}_1\) and converges in the topology \(\mathfrak{C}_2\) .

The conditions are clearly necessary, because \(\mathfrak{G}_2\) is coarser than \(\mathfrak{G}_1\) . Conversely, suppose that the conditions are satisfied, and let \(x\) be a limit point of \(\mathfrak{F}\) with respect to \(\mathfrak{G}_2\) ; we shall show that \(x\) is a limit of \(\mathfrak{F}\) with respect to \(\mathfrak{G}_1\) . Let \(V\) be a symmetric entourage of \(\mathfrak{U}_1\) which is closed in the topology \(\mathfrak{G}_2 \times \mathfrak{G}_2\) . By hypothesis, \(\mathfrak{F}\) contains a set \(M\) which is \(V\) -small; hence if \(x' \in M\) we have \(M \subset V(x')\) . But \(V(x')\) is closed in the topology \(\mathfrak{G}_{2}\) ; hence \(x\) , which lies in the closure of \(\mathbf{M}\) with respect to \(\mathfrak{G}_{2}\) , must belong to \(\mathbf{V}(x')\) . It follows that \(\mathbf{M} \subset \mathbf{V}(x)\) , and the proposition is proved.

COROLLARY. In the conditions of Proposition 7, if \(\mathfrak{U}_2\) is a uniformity of a complete space, then so is \(\mathfrak{U}_1\) .

For every Cauchy filter with respect to \(U_{1}\) is then a Cauchy filter with respect to \(U_{2}\) and therefore converges in the topology \(G_{2}\) .

Note that the hypotheses of the Corollary to Proposition 7 are satisfied when \(\mathcal{C}_1 = \mathcal{C}_2\) (§ 1, no. 2, Corollary 2 to Proposition 2).

